\color{uniblue}\section[Выбор уставок максимальной токовой защиты трансформатора]{\sloppy Выбор уставок максимальной токовой защиты \mbox{трансформатора}} 
\color{black}


\begin{enumerate}[label=\arabic{section}.\arabic{enumi}, labelsep=4pt, leftmargin=0pt, itemindent=57pt]

\item Выбор параметров срабатывания максимальной токовой защиты приводится с учетом \cite{stofsk1}.

\item Выбор параметров срабатывания первой ступени максимальной токовой защиты

\begin{enumerate}[label=\arabic{section}.\arabic{enumi}.\arabic*, labelsep=5pt, leftmargin=0em, itemindent=57pt]

\item
По условию отстройки от максимального рабочего тока трансформатора ток срабатывания защиты определяется по выражению:

\begin{equation}
\label{eq:mtz1}
I_{\textsl{с.з.МТЗ}}^I = \frac {k_{\textsl{н}} \cdot k_{\textsl{перег}}}{k_{\textsl{в}}}I_{\textsl{раб.макс.тр}},
\end{equation}
\begin{eqexpl}[15mm]
\item{$I_{\textsl{раб.макс.тр}}$} максимальное значение рабочего тока защищаемого трансформатора\footnote{~Если номинальный ток силового трансформатора больше максимального рабочего тока, то в выражение подставляется номинальный ток.}, А;
\item{$k_{\textsl{перег}}$} допустимая кратность перегрузки трансформатора в соответствии с \cite{bib:overload};
\item{$k_{\textsl{н}}$} коэффициент надёжности (отстройки), учитывающий погрешность реле и необходимый запас, принимается равным $k_{\textsl{н}}=1,2$;
\item{$k_{\textsl{в}}$} коэффициент возврата максимальных токовых органов защиты, принимается равным $k_{\textsl{в}}=0,95$.
\end{eqexpl}

\item
По условию отстройки от токов самозапуска электродвигателей ток срабатывания защиты определяется по выражению:

\begin{equation}
\label{eq:mtz2}
I_{\textsl{с.з.МТЗ}}^I = \frac {k_{\textsl{н}} \cdot k_{\textsl{сам}}}{k_{\textsl{в}}}I_{\textsl{ном.дв}\scriptscriptstyle\sum},
\end{equation}
\begin{eqexpl}[15mm]
\item{$I_{\textsl{ном.дв}\scriptscriptstyle\sum}$} сумма токов нагрузки электродвигателей, А;
\item{$k_{\textsl{сам}}$} коэффициент самозапуска нагрузки, представляющий отношение тока при самозапуске электродвигателей к предаварийному рабочему току, его значение зависит в основном от вида нагрузки (доли асинхронных двигателей, участвующих в самозапуске);
\item{$k_{\textsl{н}}$} коэффициент надёжности согласования, принимается равным $k_{\textsl{н}}=1,3$;
\item{$k_{\textsl{в}}$} коэффициент возврата максимальных токовых органов защиты, принимается равным $k_{\textsl{в}}=0,95$.
\end{eqexpl}
Если двигатели на напряжение 6-10 кВ отсутствуют, то отстройка тока срабатывания по этому условию не производится.

\item
Расчет коэффициента самозапуска электродвигателей для выражения (\ref{eq:mtz2}). Если нагрузка содержит в основном электродвигатели, то коэффициент самозапуска может быть определен по выражению:
\begin{equation}
    \label{eq:mtz3}
    k_{\textsl{сам}} = \frac {\scriptstyle\sum \textstyle I_{\textsl{пуск. двиг.}}}{\scriptstyle\sum \textstyle I_{\textsl{ном. двиг.}}},
\end{equation}
\begin{eqexpl}[15mm]
    \item{$\scriptstyle{\sum} \textstyle {I_{\textsl{пуск. двиг.}}}$} сумма пусковых токов электродвигателей, А;
    \item{$\scriptstyle{\sum} \textstyle {I_{\textsl{ном. двиг.}}}$} сумма номинальных токов электродвигателей, А.
\end{eqexpl}
В других случаях коэффициент самозапуска рассчитывается индивидуально, в зависимости от конфигурации сети, сопротивления ее элементов и соотношения мощности электродвигателей и прочей нагрузки.

\item Исходя из условий (\ref{eq:mtz1}) и (\ref{eq:mtz2}) выбирается наибольший ток срабатывания $I_{\textsl{с.з.МТЗ}}^I$.

\item
Выдержка времени выбирается на ступень селективности больше времени срабатывания самой медленно действующей защиты питающей сети
\begin{equation}
\label{eq:mtz4}
t_{\textsl{с.з.МТЗ}}^I = t_{\textsl{с.з.МТЗ}}^{\textsl{фид}}+{\Delta t},
\end{equation}
\begin{eqexpl}[15mm]
\item{$t_{\textsl{с.з.МТЗ}}^{\textsl{фид}}$} время срабатывания самой медленно действующей защиты питающей сети, с;
\item{$\Delta t$} ступень селективности (принимается равной 0,3-0,5), с.
\end{eqexpl}

\item
Чувствительность первой ступени МТЗ определяется по току наиболее неблагоприятного вида повреждения, как правило, двухфазного КЗ, в минимальном режиме работы сети на шинах одного из распределительных пунктов (РП), где он имеет наименьшее значение по выражению:

\begin{equation}
\label{eq:mtz5}
k_{\textsl{ч}} = \frac{I_{\textsl{к.мин.}}^{(2)}}{I_{\textsl{с.з. МТЗ}}^{I}},
\end{equation}
\begin{eqexpl}[15mm]
\item{$I_{\textsl{к.мин.}}^{(2)}$} ток двухфазного КЗ в минимальном режиме работы сети на шинах одного из распределительных пунктов, где он имеет наименьшее значение, А;
\item{$I_{\textsl{с.з. МТЗ}}^{I}$} ток срабатывания первой ступени МТЗ, А.
\end{eqexpl}

В соответствии с \cite{prikaz546} коэффициент чувствительности $k_{\textsl{ч}}$ должен быть $\ge 1,3$ для токовых защит трансформатора (при введенной в работу второй ступени МТЗ, обеспечивающей  $k_{\textsl{ч}}\ge 1,5$). При недостаточной чувствительности необходимо использовать пуск по напряжению со стороны НН трансформатора. 

\item
Ток срабатывания МТЗ с пуском по напряжению определяется по условию отстройки от номинального тока трансформатора по выражению:
\begin{equation}
\label{eq:mtz6}
I_{\textsl{с.з.МТЗ}}^{I} = \frac {k_{\textsl{н}}}{k_{\textsl{в}}}I_{\textsl{раб.макс.тр}},
\end{equation}
\begin{eqexpl}[15mm]
\item{$I_{\textsl{раб.макс.тр}}$} максимальное значение рабочего тока защищаемого трансформатора\footnote{~Если номинальный ток силового трансформатора больше максимального рабочего тока, то в выражение подставляется номинальный ток.}, А;
\item{$k_{\textsl{н}}$} коэффициент надёжности, учитывающий погрешность реле и необходимый запас, принимается равным $k_{\textsl{н}}=1,2$;
\item{$k_{\textsl{в}}$} коэффициент возврата максимальных токовых органов защиты, принимается равным $k_{\textsl{в}}=0,95$.
\end{eqexpl}

\item 
Минимальное рабочее напряжение на стороне НН трансформатора определяется по выражению:
\begin{equation}
\label{eq:mtz7}
U_{\textsl{с.з.НН}} = \frac {U_{\textsl{р.мин.}}}{k_{\textsl{н}} \cdot k_{\textsl{в}} \cdot k_{\textsl{сам}}},
\end{equation}
\begin{eqexpl}[15mm]
\item{$U_{\textsl{р.мин.}}$} минимальное рабочее напряжение в нагрузочном режиме, В;
\item{$k_{\textsl{н}}$} коэффициент надёжности, учитывающий погрешность реле и необходимый запас, принимается равным $k_{\textsl{н}}=1,2$;
\item{$k_{\textsl{в}}$} коэффициент возврата, принимается равным $k_{\textsl{в}}=1,05$;
\item{$k_{\textsl{сам}}$} коэффициент самозапуска нагрузки.
\end{eqexpl}

\end{enumerate}

\item Выбор параметров срабатывания второй ступени максимальной токовой защиты
\begin{enumerate}[label=\arabic{section}.\arabic{enumi}.\arabic*, labelsep=5pt, leftmargin=0em, itemindent=57pt]

\item
По условию отстройки от максимального рабочего тока трансформатора ток срабатывания защиты определяется по выражению:
\begin{equation}
\label{eq:mtz8}
I_{\textsl{с.з.МТЗ}}^{II} = \frac {k_{\textsl{н}}}{k_{\textsl{в}}}I_{\textsl{раб.макс.тр}},
\end{equation}
\begin{eqexpl}[15mm]
\item{$I_{\textsl{раб.макс.тр}}$} максимальное значение рабочего тока защищаемого трансформатора, А;
\item{$k_{\textsl{н}}$} коэффициент надёжности (отстройки), учитывающий погрешность реле и необходимый запас, принимается равным $k_{\textsl{н}}=1,2$;
\item{$k_{\textsl{в}}$} коэффициент возврата максимальных токовых органов защиты, принимается равным $k_{\textsl{в}}=0,95$.
\end{eqexpl}

\item
По условию отстройки от токов самозапуска электродвигателей ток срабатывания защиты определяется по выражению (\ref{eq:mtz7}).

\item Исходя из условий (\ref{eq:mtz8}) и (\ref{eq:mtz7}) выбирается наибольший ток срабатывания $I_{\textsl{с.з.МТЗ}}^{II}$.
\item
Выдержка времени выбирается на ступень селективности больше времени срабатывания первой ступени МТЗ:
\begin{equation}
\label{eq:mtz10}
t_{\textsl{с.з.МТЗ}}^{II} = t_{\textsl{с.з.МТЗ}}^{I}+{\Delta t},
\end{equation}
\begin{eqexpl}[15mm]
\item{$t_{\textsl{с.з.МТЗ}}^{I}$} время срабатывания первой ступени МТЗ, с;
\item{$\Delta t$} ступень селективности (принимается равной 0,3-0,5), с.
\end{eqexpl}

\item
Чувствительность второй ступени МТЗ определяется по току наиболее неблагоприятного вида повреждения, как правило, двухфазного КЗ, в минимальном режиме работы сети в конце следующего за трансформатором элемента сети по выражению:
\begin{equation}
\label{eq:mtz11}
k_{\textsl{ч}} = \frac{I_{\textsl{к.мин.}}^{(2)}}{I_{\textsl{с.з. МТЗ}}^{II}},
\end{equation}
\begin{eqexpl}[15mm]
\item{$I_{\textsl{к.мин.}}^{(2)}$} минимальный ток, протекающий на стороне ВН при двухфазном КЗ в конце последующего за трансформатором элемента сети, А;
\item{$I_{\textsl{с.з. МТЗ}}^{II}$} ток срабатывания второй ступени МТЗ, А.
\end{eqexpl}

В соответствии с \cite{prikaz546} коэффициент чувствительности $k_{\textsl{ч}}$ должен быть $\ge 1,2$ для токовых защит, выполняющих функции дальнего резервирования. При недостаточной чувствительности необходимо использовать пуск по напряжению со стороны НН трансформатора. 

\item
Ток срабатывания МТЗ с пуском по напряжению определяется по условию отстройки от номинального тока трансформатора по выражению:
\begin{equation}
\label{eq:mtz12}
I_{\textsl{с.з.МТЗ}}^{I} = \frac {k_{\textsl{н}}}{k_{\textsl{в}}}I_{\textsl{раб.макс.тр}},
\end{equation}
\begin{eqexpl}[15mm]
\item{$I_{\textsl{раб.макс.тр}}$} максимальное значение рабочего тока защищаемого трансформатора\footnote{~Если номинальный ток силового трансформатора больше максимального рабочего тока, то в выражение подставляется номинальный ток.}, А;
\item{$k_{\textsl{н}}$} коэффициент надёжности, учитывающий погрешность реле и необходимый запас, принимается равным $k_{\textsl{н}}=1,2$;
\item{$k_{\textsl{в}}$} коэффициент возврата максимальных токовых органов защиты, принимается равным $k_{\textsl{в}}=0,95$.
\end{eqexpl}

\item 
Минимальное рабочее напряжение на стороне НН трансформатора определяется по выражению (\ref{eq:mtz7}).

\end{enumerate}

\item Выбор параметров срабатывания функции блокировки при бросках токов намагничивания
\begin{enumerate}[label=\arabic{section}.\arabic{enumi}.\arabic*, labelsep=5pt, leftmargin=0em, itemindent=57pt]
    \item В целях обеспечения несрабатывания защиты при постановке трансформатора под напряжение возможна блокировка срабатывания ступеней МТЗ от функции БНТ. БНТ контролирует превышение отношения тока второй гармоники к первой заданного значения уставки. Значение уставки рекомендуется принимать равным 15\%.
    \item Для исключения возможности работы функции БНТ при больших кратностях тока, в функции БНТ реализована блокировка при превышении измеряемого тока заданной уставки. Рекомендуется принимать значение тока срабатывания блокировки равным $I_{\textsl{ср}}=7,5 \cdot I_{\textsl{ном.тр.}}$.
\end{enumerate}

\end{enumerate}